\documentclass[12pt,a4paper]{article}
\usepackage[T1]{fontenc}
\usepackage[margin=1in]{geometry}
\usepackage{enumitem}

\begin{document}

\begin{center}
  {\Large\bfseries Worked Solutions: The Lattice Multiplication Method -- Answer Key}
\end{center}

\section*{Q1}
(a) \(52\times34=1768\)\quad
(b) \(68\times47=3196\)

\section*{Q2}
(a) \(34\times25=850\)\quad
(b) \(243\times17=4131\)

\section*{Q3}
\(1.7\times4.3=7.31\)

\section*{Q4}
(a) \(0.6\times4.5=2.7\)\quad
(b) \(12.4\times0.3=3.72\)

\section*{Q5}
(a) \(6\times45=270\)\quad
(b) \(0.06\times4.5=0.27\)\quad
(c) \(60\times0.45=27\).
The digits \(27\) stay the same; only the place value changes.

\section*{Q6}
(a) \(24\times3.65=\pounds87.60\)\quad
(b) \(0.35\times16=5.60\), so \(5.6\) litres are lost.

\section*{Q7}
(a) \(1.45\times38=\pounds55.10\)\quad
(b) \(4.8\times2.5=12\ \mathrm{cm}^2\)

\section*{Q8}
(a) \(73\times48=3504\)\quad
(b) \(5.2\times3.6=18.72\)\quad
(c) \(0.9\times0.7=0.63\)\quad
(d) \(126\times34=4284\)\quad
(e) \(8.4\times0.25=2.1\)\quad
(f) \(1.06\times3.5=3.71\)

\section*{Q9}
(a) \(40\times6=240\) and \(7\times30=210\); both are whole numbers of tens, so both belong to the tens column.
(b) Cells on the same lattice diagonal contribute to the same place value, so their digits are added; totals of \(10\) or more carry into the next diagonal.

\section*{Q10}
(a) \(0.48\times0.25\) must be less than \(1\), so \(1200\) is impossible.
(b) \(0.48\times0.25=0.12\)

\section*{Q11}
\(3\times4=12\) single-digit multiplications.

\section*{Q12}
Example: \(0.4\times0.9=0.36\).

\section*{Q13}
Zainab is partly right: no carrying is needed while filling the cells, but carrying is used when adding along the diagonals.

\end{document}